% Appendix E: the attacks and defenses of Section 5 (split from Appendix D, session 106). Session 108: mirrors §5's
% subsections (memory: appendix mirrors body); no subsection for 5.2, which needs none. Rule for added defenses
% (docs/evaluation-note.md, decision 5): a row in the comparison table (tab:eval:detectors) and the per-attack table, a paragraph in E.2 on how it was
% built and ported.

This appendix gives the details of the attacks and defenses that \Cref{sec:evaluation} runs, in the order of its
subsections.

\subsection{Setup}
\label{sec:attacks:setup}

\paragraph{The toolkit attacks.} Each runs secml-malware's genetic algorithm at the budget of the GAMMA
paper~\cite{demetrio2021functionality}: a population of 10 over 51 iterations, with 75 benign donor files. The five
byte-level attacks were built for models that read the raw bytes, and against our detectors they rarely evade
(\Cref{tab:ember:per-attack}). Against secml-malware's pretrained MalConv~\cite{raff2018malware}, the model they were
built for, the same code and budget evade between 5 and 16 of 20 malware files, so neither the code nor the budget is
why they rarely evade ours. The released GAMMA attacks write their sections and imports through the LIEF library~\cite{lief}, which
rebuilds the file's layout and moves content that no operation moves. We keep GAMMA's search and write the edits it
chooses with our own code, growing the headers first where they lack room for the new section headers; where no growth
makes room, the attack does not apply to the file.
% docs/scaleup-protocol.md item (22): diag/s105/malconv_control.py (16, 14, 9, 6, 5 of 20) and toolkit_reach.py.

\paragraph{The adaptive attack.} The adaptive attack of \Cref{sec:eval:setup} is a beam search. At each step it applies
random candidate edits to the three lowest-scoring files found so far, writes each candidate with our own code and
extracts its features, and keeps the three lowest-scoring files; after three steps without improvement it restarts
from the original file. When it adds or renames a section, it picks a name whose hash lands in a bucket the detector's
trees weigh toward benign. When it adds a section or imports and the headers lack room for them, it grows the headers
first, as the GAMMA attacks do.
% pipeline/5-attacks/adaptive (adaptive.py, --score real at the scale-up; decision 12 (iii)); target_buckets reads the
% booster; hgr added in session 105 (item (22)).

\paragraph{MAB-Malware.} We run the authors' code unchanged except for how it hands files to the detector. Upstream,
the attack, its minimizer and the detector run as three threads that pass files through folders on disk; we call the
same steps in the same order from one thread and score each file when it is handed over, and we checked on test files
that this produces the same actions, outputs and bandit state as the original hand-over. One run attacks all the
malware a detector detects, so the bandit learns from one detector at a time. Zainab et al.~\cite{zainab2026breaking}
report that MAB-Malware evades their EMBER~2024 detector on 83\% of the malware, at a 0.1\% false-positive rate. Their
detector and malware differ from ours, so the two rates describe different setups: at our 1\% rate, it evades the
unconstrained detector of \Cref{tab:eval:detectors} on 10 of the 47 malware it detects, and the same detector trained
on the full training set on none of 49.
% Old round's values (5%/512 KB cap); the new round (4 MB rule, session 105) replaces them, and adds MAB-Malware at
% Zainab et al.'s setting (0.1% FPR, no byte limit) against the unconstrained detector (item (8)).
% Zainab et al., Table I (MAB-Malware on EMBER-24, 83.40%) and Table IX (EMBER-24 at 0.1% FPR); no byte limit stated.
% scripts/mab_attack.py (MABChannel, run_sync); tests/test_mab_attack.py; third_party/mab_malware (upstream 91a518a).

\paragraph{Each attack separately.} \Cref{tab:ember:per-attack} gives, for every detector of
\Cref{tab:eval:detectors}, the malware each attack evades.

\begin{table*}[t]
  \caption{Malware evaded by each attack, per detector, grouped as in \Cref{tab:eval:detectors}. Toolkit attacks as in
  \Cref{tab:instances:catalogue}.}
  \label{tab:ember:per-attack}
  \small
  \begin{tabular*}{\textwidth}{@{}l@{\extracolsep{\fill}}rrrrrrrrrrr@{}}
    \toprule
    Detector & detected & padding & \multicolumn{2}{c}{GAMMA} & GAMMA & content & DOS & header & format & adaptive
      & MAB- \\
    & & & sections & imports & shift & shift & ext. & fields & exploit & & Malware \\
    \midrule
    Unconstrained        & 47 & 1 & 6 & 24 & 7 & 0 & 0 & 9 & 0 & 47 & 10 \\
    \midrule
    \multicolumn{12}{@{}l}{\emph{Read by hand}} \\
    Expert set (additions) & 8  & 0 & 0 & 0  & 0 & 0 & 0 & 0 & 0 & 2  & 4 \\
    \midrule
    \multicolumn{12}{@{}l}{\emph{Learned from observed edits}} \\
    PV                   & 9  & 0 & 0 & 0  & 0 & 0 & 0 & 0 & 0 & 4  & 7 \\
    ERDALT               & 29 & 0 & 0 & 24 & 0 & 0 & 0 & 0 & 0 & 22 & 6 \\
    Adversarial training & 50 & 0 & 1 & 10 & 4 & 2 & 0 & 2 & 0 & 23 & 4 \\
    Iterated adv.\ training & \synthetic{50} & \synthetic{0} & \synthetic{1} & \synthetic{5} & \synthetic{2}
      & \synthetic{1} & \synthetic{0} & \synthetic{1} & \synthetic{0} & \synthetic{48} & \synthetic{4} \\
    Trial                & 44 & \synthetic{0} & \synthetic{0} & \synthetic{0} & \synthetic{0} & \synthetic{0}
      & \synthetic{0} & \synthetic{0} & \synthetic{0} & 26 & 17 \\
    \midrule
    \multicolumn{12}{@{}l}{\emph{Derived from the extractor}} \\
    Monotone, features   & 39 & \synthetic{0} & \synthetic{0} & \synthetic{0}  & \synthetic{0} & \synthetic{0} & \synthetic{0} & \synthetic{0} & \synthetic{0} & 0  & \synthetic{0} \\
    Monotone, sums       & 40 & \synthetic{0} & \synthetic{0} & \synthetic{0}  & \synthetic{0} & \synthetic{0} & \synthetic{0} & \synthetic{0} & \synthetic{0} & 0  & \synthetic{0} \\
    Monotone, sums (additions) & 43 & 0 & 0 & 0  & 0 & 0 & 0 & 12 & 0 & --- & 1 \\
    \bottomrule
  \end{tabular*}
\end{table*}
% Session 108: regrouped as the comparison table (tab:eval:detectors); unconstrained is s10-ref (run 143 adaptive with
% the entry write; toolkit a1/toolkit-v2; MAB run 139). PV, ERDALT, adversarial training: a_varied/toolkit-*/table.json,
% a_varied/adaptive-*-best, run 139. Expert set: a1/expert/toolkit, run 120, run 139. Iterated: SYNTHETIC (no pilot run;
% predicted near adversarial training's toolkit cells, the adaptive attack through `ent`). Trial: run 144 (beam at
% 2,000 candidates, MAB-Malware); its toolkit cells SYNTHETIC (no toolkit run on it; predicted 0: its table keeps the
% import buckets nondecreasing, which GAMMA imports raises). Monotone rows: run 143 adaptive measured; toolkit and
% MAB-Malware cells predicted 0 on the -loose detectors. Additions: toolkit a1/toolkit-v2 (12 header fields); no
% adaptive run (---); MAB run 139. SCALE-UP: one protocol and one detected count per row (R11); 6-results.

\subsection{What Observed Edits Miss}
\label{sec:attacks:observed}

\paragraph{PV.} PV selects among EMBER's features, not the sums. A feature counts as lowered when some output lowers
it by more than $10^{-6}$, which only guards against rounding; the lowered features are left out of the detector.
% a_pv_hole.py rule (kept = min observed delta >= -tol, monotone +1); at the scale-up pv_hole.py --exclude in
% 4-detectors/learned, selected among EMBER's features (the re-encoded sums dropped).

\paragraph{ERDALT.} We follow Gimenez et al.'s released code with one change: the output layer's weights are kept
nonnegative, like those of the layers below it, so that the network is monotone as their Proposition~3 requires; the
released code leaves the output layer unconstrained. Their code sizes several settings to their data, 24,000 training
files and 1,033 adversarial feature changes, against our 1.76 million and \synthetic{about 80,000}. We rescale each so
that it has the same effect: the penalty on the layer is weighted by 1,033 over the number of changes a model sees,
so that it weighs as much per step as in their training; the schedule covers their number of training steps, 68 of our
epochs, with the learning rate's patience shortened in proportion; the features are scaled by their quantiles; and
their selection criterion, fewer than 5 violated entries over all changes, becomes the same rate per change, fewer than
381. Their batch size of 16 becomes 256. No epoch meets the rescaled criterion, so we keep the one with the fewest
violations.
% docs/erdalt-diagnosis-plan.md (sessions 82-100); src/provexp/models/erdalt.py (project_monotone); docs/upstream-
% issues.md, item 8. Session 100: validation AUC 0.978, test 0.981, 35,594 violations against 381 (their literal w2:
% 769 violations at 0.914), kept as is (user). SYNTHETIC: the material's count (78,715 vectors in the old round).

\paragraph{Adversarial training.} The adversarially trained detector is the unconstrained detector retrained with the
outputs of the toolkit attacks added to its malware. Iterated adversarial training starts from it. Each round runs
the adaptive attack, with every operation except the entry write, on 1,000 training malware not used before, adds the
evasive outputs to the malware, and retrains; we run two rounds.
% 4-detectors/adversarial-training (material), learned (advtrain), iterated (two rounds; scaleup-protocol.md (7),
% (11): starting point chosen by Claude, the adversarially trained detector).

\paragraph{Trial.} The trial table learns from the edits of the checks of \Cref{sec:eval:setup} on the 2,000 files of
the evaluation set and from a second pass over the same files: every section renamed, the section slack, the DOS
header and stub, and the loader-ignored header fields all overwritten at once, with zeros, random bytes, and printable
text, and imports added from 100 benign training files. A feature that some edit lowered is dropped, one that some
edit raised and none lowered is monotone increasing, and one that no edit moved is left free.
% 3-directions/trial (trial_probe.py, trial_table.py; donor set 0). Trial-toolkit (pilot only, run 144): the trial
% table widened with the toolkit attacks' outputs.

\paragraph{Leave-one-out.} \Cref{tab:ember:loo} gives the folds of \Cref{sec:eval:observed}. Each fold builds PV and
ERDALT from the outputs of seven toolkit attacks and runs the eighth on the test malware the fold's detector detects,
as in \Cref{sec:eval:setup}.

\begin{table}[t]
  \caption{Leave-one-out over the toolkit attacks: test malware each fold's detector detects, and how many of them the
  held-out attack evades.}
  \label{tab:ember:loo}
  \small
  \begin{tabular}{@{}lrrrr@{}}
    \toprule
    Held-out attack & \multicolumn{2}{c}{PV} & \multicolumn{2}{c}{ERDALT} \\
    & detected & evaded & detected & evaded \\
    \midrule
    GAMMA sections & 19 & 0 & 27 & 1 \\
    GAMMA imports & 9 & 0 & 29 & 25 \\
    GAMMA shift & 9 & 0 & 30 & 0 \\
    content shift & 9 & 0 & 30 & 1 \\
    DOS extension & 9 & 0 & 28 & 0 \\
    header fields & 9 & 0 & 29 & 2 \\
    padding & 9 & 0 & 28 & 1 \\
    format exploit & 8 & 0 & 29 & 0 \\
    \bottomrule
  \end{tabular}
\end{table}
% Run 135: artifacts/experiments/a_loo_varied/{pv,erdalt}/<attack>/table.json. Attack names as in Table 2.

\subsection{What Robustness Costs}
\label{sec:attacks:cost}

\Cref{tab:ember:price} gives the detection of the monotone detector on the sums with each operation left out of the
threat model in turn, the numbers \Cref{sec:eval:detectors} draws on.

\begin{table}[t]
  \caption{The monotone detector on the sums, measured on the threat model without the entry write, with one operation
  at a time left out of it. TPR at each false-positive rate.}
  \label{tab:ember:price}
  \small
  \begin{tabular*}{\linewidth}{@{}l@{\extracolsep{\fill}}rr@{}}
    \toprule
    Detector & 1\% & 0.1\% \\
    \midrule
    Unconstrained & 93.5 & 84.4 \\
    Two-sided features dropped, no constraints & 93.3 & 79.2 \\
    Threat model without the entry write & 89.6 & 74.7 \\
    \midrule
    \multicolumn{3}{@{}l}{\emph{Leaving out}} \\
    add imports & 91.7 & 78.8 \\
    the three header overwrites & 91.0 & 79.2 \\
    \quad overwrite header fields & 90.8 & 77.0 \\
    \quad overwrite the checksum & 90.3 & 76.1 \\
    \quad overwrite the DOS header and stub & 89.4 & 75.7 \\
    remove the debug entry & 90.0 & 75.2 \\
    remove the certificate entry & 90.0 & 75.0 \\
    rename a section & 89.7 & 76.0 \\
    gap before the PE header & 89.6 & 77.1 \\
    append bytes at the end & 89.3 & 71.7 \\
    \bottomrule
  \end{tabular*}
\end{table}
% MEASURED (2026-11-04): run/126_price_list.sh, artifacts/experiments/a_price_list/table.json (before the entry write;
% ext = gap before the PE header). Not measured: add a section, gap before the first section, grow the headers,
% section slack (the price list's operations then). PROTOCOL: the scale-up's 4-detectors/price-list re-runs it on the
% widened table, every operation, including the entry write; this table becomes the full threat model's.
